\section{Assay resolution and recipient coexpression}
\label{sec:resolution}
This section gives the theory and retrospective analyses added after the
original fixed-margin benchmark. All resolution comparisons hold the source
interaction fixed.

\subsection{Sharp coarse prediction bound}\label{app:proof}
Let $x,y$ range over finite assay-profile supports, let $S(x,y)$ be finite,
and consider positive joint distributions
\begin{equation}
 q(x,y)=\exp\{S(x,y)+a(x)+b(y)\}.
 \label{eq:resolution_joint}
\end{equation}
The binary queries $U=f(x)$ and $V=g(y)$ each have two nonempty fibers and
known frequencies $m_U,m_V\in(0,1)$. For every rectangle crossing their four
fibers, define
\[
 d=S(x_0,y_0)+S(x_1,y_1)-S(x_0,y_1)-S(x_1,y_0),
\]
where $f(x_0)=g(y_0)=0$ and $f(x_1)=g(y_1)=1$. Write $L=\min d$ and
$H=\max d$. The closure of compatible coarse log odds is $[L,H]$.
The following proof also gives the minimax prediction under these two
frequency constraints.
For the known margins $m_U,m_V$, write the compatible binary table as
\[
 P(t)=\begin{pmatrix}1-m_U-m_V+t&m_V-t\\m_U-t&t\end{pmatrix}.
\]
Let $t(\theta)$ be the unique interior cell probability giving log odds
$\theta$, and put $l=t(L)$ and $h=t(H)$.

Every fine distribution in Eq.~\eqref{eq:resolution_joint} satisfies
\[
 q(x_0,y_0)q(x_1,y_1)=e^d q(x_0,y_1)q(x_1,y_0).
\]
Summing over rectangles in the four coarse blocks expresses the coarse odds
ratio as a positively weighted average of $e^d$. Its log therefore belongs to
$[L,H]$. Concentrate the fine RNA masses $1-m_U,m_U$ on an extremizing $x_0,x_1$
and protein masses $1-m_V,m_V$ on $y_0,y_1$. Their two-by-two reconstruction attains
the corresponding endpoint. Positive fine distributions approximate these
concentrated distributions while preserving the coarse margins. The positive
fine-marginal domain is connected, and its entropic reconstruction is
continuous, so the closure of attainable coarse tables is exactly
$\{P(t):l\leq t\leq h\}$.

For a fixed prediction, KL loss is convex in the true table, so its maximum on
this segment occurs at an endpoint. Moving a predicted table along the segment
increases its loss to one endpoint and decreases its loss to the other. There
is a unique equalizing table $P_*$, which is a mixture of the endpoints. For
$l<h$, define entropy as $\mathcal H(P)=-\sum P\log P$. The minimax log odds are
\[
 \theta_* = \frac{\mathcal H(P(l))-\mathcal H(P(h))}{h-l},
 \qquad P_*=P(t(\theta_*)).
\]
For $l=h$, set $P_*=P(l)$.

The equalizing construction is an information-radius minimax rule
\cite{vanerven}. To prove optimality over all predictions, including those with different
margins, choose the weight $w$ such that $P_*=wP(l)+(1-w)P(h)$. For any
prediction $R$,
\begin{align*}
 &w\KL(P(l)\Vert R)+(1-w)\KL(P(h)\Vert R)\\
 &\quad=R_*+\KL(P_*\Vert R),
\end{align*}
where $R_*=\KL(P(l)\Vert P_*)=\KL(P(h)\Vert P_*)$. The worst endpoint loss
cannot be smaller than this weighted average. Equality is attained uniquely
by $P_*$. All compatible populations induce the same independent unpaired
Bernoulli observation laws with frequencies $m_U,m_V$. Convexity in the predicted
distribution makes randomization unhelpful, so $R_*$ is also a lower bound
for every sample-based predictor observing only these unpaired binary states.

For the three-marker example in the main text, take the RNA distribution
$r(z_1,z_2)=(1+\rho z_1z_2)/4$, with $|\rho|<1$. The joint distribution
\[
 q(z_1,z_2,w)=r(z_1,z_2)\frac{e^{Jz_2w}}{2\cosh J}
\]
has all three individual means zero and the required fine interaction.
Conditional expectation gives
$\mathbb E(W\mid Z_1,Z_2)=Z_2\tanh J$ and hence
$\mathbb E(Z_1W)=\rho\tanh J$. All probabilities are positive, so the example does
not depend on structural zeros.

\subsection{Response to a transferred interaction}\label{app:response}
The expansion below specializes the fixed-marginal response identities of
Sei and Yano \cite{sei}. Fix positive marginal masses $r(x)$ and $c(y)$ on finite supports. Write
$\bar x=x-\mathbb E_r X$, $\bar y=y-\mathbb E_c Y$, and
$h=\bar x^\top B\bar y$. Centering changes the interaction only by additive
assay functions. The reconstruction can therefore be written
\[
 q_t(x,y)=r(x)c(y)\exp\{t h(x,y)+\alpha_t(x)+\beta_t(y)\}.
\]
At $t=0$, the adjusting functions vanish. After fixing their additive gauge,
the finite marginal equations have a nonsingular linearization at independence,
so their solution is analytic near zero. Differentiating the marginal
constraints once gives zero conditional means for
$\partial_t\log q_t$. Centering $h$ in both arguments therefore gives
$\partial_t\log q_t|_{t=0}=h$. Differentiating twice gives
\begin{align*}
 \left.\frac{\partial_t^2q_t}{rc}\right|_{t=0}
 &=h^2-\mathbb E_c(h^2\mid x)\\
 &\quad-\mathbb E_r(h^2\mid y)+\mathbb E_{rc}h^2.
\end{align*}
When multiplied by $\bar x_i\bar y_j$ and averaged under $rc$, the last
three terms vanish. Define the third central moment tensors by
$T^x_{iac}=\mathbb E_r(\bar X_i\bar X_a\bar X_c)$ and
$T^y_{jbd}=\mathbb E_c(\bar Y_j\bar Y_b\bar Y_d)$. It follows that
\begin{align*}
 &\operatorname{Cov}_{q_t}(X_i,Y_j)=t(\Sigma_x B\Sigma_y)_{ij}\\*
 &\quad+\frac{t^2}{2}\sum_{a,b,c,d}
       B_{ab}B_{cd}T^x_{iac}T^y_{jbd}+O(t^3).
\end{align*}
For an explicit example, take three spins in each assay and let
\[
 r(x)=\frac{1+\rho_x x_1x_2x_3}{8},\qquad
 c(y)=\frac{1+\rho_y y_1y_2y_3}{8},
\]
with $x_i,y_i\in\{-1,1\}$ and $|\rho_x|,|\rho_y|<1$. Every individual mean
is zero and each within-assay covariance matrix is the identity. For diagonal
$B=\operatorname{diag}(b_1,b_2,b_3)$,
\[
 \mathbb E_{q_t}(X_1Y_1)
 =t b_1+t^2\rho_x\rho_y b_2b_3+O(t^3).
\]
The pairwise maximum-entropy replacement is uniform in each assay and loses
the second term. Positive distributions with identical means and covariances
can therefore give different RNA--protein associations under the same
bilinear interaction.

\subsubsection{Query discrepancy from covariance mismatch}\label{app:query}
Consider recipient and pooled reconstructions with the same binary marker
frequencies $m,n\in(0,1)$ and interaction $t x^\top B y$. Let $k_R$ and $k_P$
be their respective queried entries of $\Sigma_x B\Sigma_y$. Both query tables
equal the independent table $P_0$ at $t=0$, and
\[
 P_R(t)-P_P(t)=t(k_R-k_P)V+O(t^2),
\]
where $V=\left(\begin{smallmatrix}1&-1\\-1&1\end{smallmatrix}\right)$.
Taylor expansion of KL divergence about the positive table $P_0$ gives
\begin{align*}
 &2\KL(P_R(t)\Vert P_P(t))\\*
 &\quad=t^2(k_R-k_P)^2\sum_{u,v}\frac{V_{uv}^2}{(P_0)_{uv}}+O(t^3)\\*
 &\quad=\frac{t^2(k_R-k_P)^2}{m(1-m)n(1-n)}+O(t^3).
\end{align*}
The coefficient is the response score used to rank queries. If the recipient
model describes the true joint law, it is also the leading excess deviance
from substituting pooled coexpression. The empirical comparison tests the
ranking at the fitted interaction strength without imposing that equality.



\subsection{Matched resolution and covariance analyses}

The three recipient cohorts contain 24 Cambridge donors, 56 Newcastle donors,
and 15 T-ALL patients. The same nine RNA and nine protein markers give 81
ordered queries. All analyses retain the original deterministic 512-cell
sample, donor-wide protein ranks, 256 adaptation cells, 256 scoring cells,
and 12-donor source interaction with pairing-contrast penalty $10^{-4}$.
The paired source fit is the only source of cross-assay interaction
information. Recipient adaptation uses separate-assay patterns and the
scoring sample's individual marker frequencies, not its pairings.

The means-only reconstruction constrains all 18 individual frequencies.
Each one-pattern arm additionally constrains that assay's full binary-profile
histogram. The both-pattern arm constrains both histograms. Histograms receive
four source-prior pseudocells and exponential tilts to the scoring-cell
marker frequencies. Constant markers are restricted to their observed states.
Matrix scaling fits complete pattern constraints, and moment matching fits
individual-frequency constraints. These are reconstructions from the same
interaction, not separately trained interaction models.

The covariance arm replaces each smoothed and tilted histogram with its
maximum-entropy distribution subject to the same first and second moments.
For binary profiles its log density contains unary and pairwise terms.
This replacement can induce higher-order correlations; it fixes pairwise
moments rather than forcing all higher-order moments to vanish.
The retained-gain fraction is
\[
 \frac{L_{\mathrm{means}}-L_{\mathrm{covariance}}}
 {L_{\mathrm{means}}-L_{\mathrm{patterns}}},
\]
where each $L$ is the cohort's mean donor deviance. It is a descriptive loss
decomposition, not an estimate of a mediated biological effect.

\subsection{Recipient-specific and pooled coexpression}

The source pool uses the 12 source donors' adaptation halves: 3,072 cells
per assay. For each query person, the cohort pool uses every other recipient's
adaptation half and excludes the query person. It contains 5,888 cells in
Cambridge, 14,080 in Newcastle, and 3,584 in T-ALL. Both controls use four prior
pseudocells per contributing person, preserving the prior fraction, then
undergo the same tilting and maximum-entropy reconstruction at the query
person's marker frequencies. The cohort pool uses recipient-cohort data but
no scoring-cell patterns. These controls vary the origin of covariance
without changing the interaction, marker frequencies, or scoring tables.

\subsection{Query ranking and uncertainty}

For RNA marker $i$ and protein marker $j$, define
\[
 s_{ij}=
 \frac{\{(\Sigma_x B\Sigma_y)^{\mathrm{recipient}}_{ij}
       -(\Sigma_x B\Sigma_y)^{\mathrm{pool}}_{ij}\}^{2}}
 {m_i(1-m_i)n_j(1-n_j)}.
\]
The source-strength comparator is
$B_{ij}^{2}m_i(1-m_i)n_j(1-n_j)$.
Scores are zero for constant-marker queries. Both rankings use the fixed
marker order to break ties. Pair-specific gain is the cohort-pool deviance
minus recipient-covariance deviance, using unchanged predictions.
The top 27, bottom 27, and source-ranked top 27 were specified before
computing the pair-specific losses. Cumulative gain curves sum gains over
people and selected queries, then divide by the sum over all 81 queries;
they do not average person-specific percentages.

All loss endpoints average twice KL over every marker pair, including
constants, then equally over people. Each comparison uses 20,000 paired
person-level bootstrap draws. The four analysis families use seeds
906091--906094 and contain, respectively, nine pattern-resolution contrasts,
six covariance-replacement contrasts, six recipient-versus-pool contrasts,
and six query-ranking contrasts. For a family of $M$ comparisons, the
Bonferroni percentile endpoints are $0.025/M$ and $1-0.025/M$.
These intervals condition on the fixed source fit and fitted pools.

\begin{table}[htbp]
\caption{Query-ranking comparisons. Values are mean deviance-gain differences
per selected pair. Intervals resample people and correct across six contrasts.}
\centering\small
\begin{tabularx}{\textwidth}{@{}lYrrY@{}}
\toprule
Cohort & Comparison & Difference & Favorable people & Bonferroni interval\\
\midrule
Cambridge & Response top vs bottom third & 0.001711 & 15/24 & $-0.000126$ to 0.003883\\
Newcastle & Response top vs bottom third & 0.004983 & 44/56 & 0.002657 to 0.007787\\
T-ALL & Response top vs bottom third & 0.024259 & 14/15 & 0.014180 to 0.037793\\
Cambridge & Response vs source-ranked top third & 0.000038 & 11/24 & $-0.001043$ to 0.001154\\
Newcastle & Response vs source-ranked top third & 0.001642 & 38/56 & 0.000612 to 0.002920\\
T-ALL & Response vs source-ranked top third & 0.013881 & 14/15 & 0.007619 to 0.021405\\
\bottomrule
\end{tabularx}
\end{table}

\subsection{Relation to full-profile interaction fitting}

The earlier T-ALL comparison used all 15 patients, the same permitted
recipient information, and additional protein-abundance coordinates.
A common recipient adapter compared the no-prior CHAMPOLLION inverse-transport
objective \cite{champollion} with donor-wise joint fitting, query-targeted
fitting, and a multinomial comparator. Mean losses were 0.008441, 0.008476,
0.008602, and 0.010630, respectively. Query-targeted fitting improved on the
multinomial comparator by 19.07\% (95\% interval, 13.05--24.36\%); its relative
improvement over the CHAMPOLLION objective was $-1.91\%$
($-6.26$ to 2.25\%). This comparison did not meet its all-control superiority
criterion. It neither benchmarks the complete CHAMPOLLION workflow nor
establishes a better interaction learner. The subsequent binary resolution
analyses hold one fitted interaction constant.

\subsection{Computational verification}

Separate implementations reconstructed all 380 initial resolution predictions,
190 covariance-preserving assay distributions and their 95 joint predictions,
and 190 pooled-control predictions. Maximum absolute differences in predicted
probabilities were $5.04\times10^{-11}$, $1.14\times10^{-10}$, and
$1.27\times10^{-10}$, respectively. Verification recounted the scoring tables,
checked exact marker margins, and checked source membership and exclusion of
the query person from each cohort pool. Independently calculated query
scores differed by at most $1.47\times10^{-10}$; pair-specific gain and
bootstrap-summary differences were below $4\times10^{-16}$.
Fourteen tests cover marginal constraints, reconstruction equivalence,
covariance replacement, pool construction, and query scoring.

Analysis plans preceded these computations but followed examination of earlier
outcomes; the analyses are retrospective. The original public v2.0.4 benchmark
is a separate release and does not contain the subsequent resolution,
covariance, pooled-control, or query-ranking analyses.
